% =====================================================================
%  Lecture 2026-09-10
%
%  CORRECTED SOURCE ERRORS
%   [p1] Hölder stated for 1 <= p <= infinity, proved only for
%        1 < p < infinity  ->  hypothesis 1 < p < infinity, endpoint
%        pair (1, infinity) stated and proved separately
%   [p1] Hölder statement indexed the left sum by i and one right factor
%        by j  ->  index j throughout
%   [p1] final Hölder chain "= 1 = ||x||_p ||y||_q" held only under the
%        normalisation  ->  scaling step written out, equality now true
%        as stated
%   [p1] (p-1)q = p rearranged to q = (p-1)/p  ->  q = p/(p-1)
%   [p2] L_p defined with "Lebesgue integrable"  ->  measurable
%   [p2] general measure-space remark asserted Banach for
%        1 <= p <= infinity  ->  1 <= p < infinity, matching the norm
%        formula the source gives
%   [p2] completion called unique  ->  unique up to isometric isomorphism
%   [p2] lemma proof "Goal: ||sum a_i x_i|| > C"  ->  >= C
%   [p2] exercise 1 chain used a strict inequality between the two tail
%        sums  ->  non-strict
%
%  ADDITIONS
%   [p1] definitions of metric space, Cauchy sequence and completeness,
%        named but not stated in the recap
%   [p1] definition of a norm, implied by "checking the normed
%        properties" but never stated
%   [p1] connectedness hypothesis in the shortest-path example
%   [p1] Young's inequality: cases a = 0 and b = 0 separated, positivity
%        of u_1 and u_2 noted, concavity of log justified by its second
%        derivative
%   [p1] Minkowski: the degenerate case sum |x_j + y_j|^p = 0 excluded
%        before dividing, and 1 - 1/q = 1/p made explicit
%   [p2] L_p[0,1] defined as the completion of (C[0,1], ||.||_p); the
%        source's set description kept as a remark
%   [p2] exercise 1: the existence of N justified by convergence of the
%        series
%   [p2] exercise 2: the bound on the second integral attributed to
%        mu([0,1]) = 1, and a pointer to the representation remark
%   [p2] definition of equivalent norms, named in "Recall" but not stated
%   [p2] lemma proof: reduction to sum |a_i| = 1 by homogeneity,
%        Bolzano-Weierstrass named, continuity of the norm and of the
%        coefficient map named
%   [p3] max over i = 1, ..., n quantified; transitivity of equivalence
%        written out
%
%  REPORTED BUT UNMODIFIED GAPS
%   [p1] concavity of log is used, not proved
%   [p2] incompleteness of (C[0,1], ||.||_p) asserted without a
%        counterexample
%   [p2] that C[0,1] is dense in L_p[0,1] is implicit in the definition
%        by completion and is not discussed
%   [p2] Hölder and Minkowski for sequences and for functions stated
%        without proof, as in the source and the syllabus
%   [p2] completion proposition stated without proof
%   [p2] exercise 2 argues on representatives; that the induced map
%        L_{p'}[0,1] -> L_p[0,1] is injective is not established
%   [p3] ||.||_1 asserted to be a norm without verification
%
%  OMITTED
%   [p1] the Wasserstein distance passage in full, with its two figures
%   [p2] boxed homework note on L_p(R) and L_{p'}(R), one word illegible
%   [p3] closing exercise on the equivalence of ||.||_1 and ||.||_2 on
%        ell_1, unsolved in the source
%   [p1, p2] doodles unconnected to the material
% =====================================================================

% [p1]
\section{Metric spaces}

\begin{definition}[Metric space]\label{def:metric}
Let $X$ be a set. A function $d \colon X \times X \to \R$ is a \emph{metric}
if for all $x, y, z \in X$ it satisfies $d(x,y) \ge 0$ with equality exactly
when $x = y$, the symmetry $d(x,y) = d(y,x)$, and the triangle inequality
$d(x,z) \le d(x,y) + d(y,z)$. The pair $(X,d)$ is a \emph{metric space}.
\end{definition}

\begin{definition}[Cauchy sequence]\label{def:cauchy}
A sequence $(x_n)$ in a metric space $(X,d)$ is \emph{Cauchy} if for every
$\eps > 0$ there is an $N \in \N$ such that $d(x_n, x_m) < \eps$ for all
$n, m > N$.
\end{definition}

\begin{definition}[Completeness]\label{def:complete}
A metric space $(X,d)$ is \emph{complete} if every Cauchy sequence in $X$
converges to a limit lying in $X$.
\end{definition}

\begin{example}[Shortest-path metric]\label{ex:graph}
Let $G = (V,E)$ be a connected graph and let $d(v_1, v_2)$ be the length of a
shortest path between $v_1$ and $v_2$. Then $d$ is a metric on $V$.
Connectedness is what makes $d$ finite.
\end{example}

\begin{proposition}[Completion]\label{prop:completion}
For every metric space $(X,d)$ there is a complete metric space
$(\cl{X}, \cl{d})$ together with a map $i \colon X \to \cl{X}$ such that
$\cl{d}(i(x), i(y)) = d(x,y)$ for all $x, y \in X$ and $i(X)$ is dense in
$\cl{X}$. It is unique up to isometric isomorphism: any two such spaces are
related by a bijective isometry carrying one copy of $X$ onto the other.
$(\cl{X}, \cl{d})$ is the \emph{completion} of $(X,d)$.
\end{proposition}

\begin{example}
The space $(\Q, \abs{\cdot})$ is not complete; its completion is
$(\R, \abs{\cdot})$.
\end{example}

\begin{remark}
It is common to suppress $i$ and assume outright that $X \subset \cl{X}$, in
which case the last condition reads: $X$ is dense in $\cl{X}$.
\end{remark}

% [p1]
\section{Normed spaces and the \texorpdfstring{$p$}{p}-norms}

\begin{definition}[Norm]\label{def:norm}
Let $X$ be a vector space over $\K$, where $\K$ is $\R$ or $\C$. A function
$\norm{\cdot} \colon X \to \R$ is a \emph{norm} if for all $x, y \in X$ and
$\lambda \in \K$ it satisfies $\norm{x} \ge 0$ with equality exactly when
$x = 0$, the homogeneity $\norm{\lambda x} = \abs{\lambda}\norm{x}$, and the
triangle inequality $\norm{x+y} \le \norm{x} + \norm{y}$. Then
$(X, \norm{\cdot})$ is a \emph{normed space}, and $d(x,y) = \norm{x-y}$ is a
metric on $X$.
\end{definition}

\begin{example}\label{ex:spaces}
\leavevmode
\begin{enumerate}
  \item $X = \K^n$ with the $p$-norm $\norm{x}_p = (\sum_{j=1}^n
    \abs{x_j}^p)^{1/p}$ for $1 \le p < \infty$, or with
    $\norm{x}_\infty = \max_{1 \le j \le n} \abs{x_j}$.
  \item $X = \Cont{0,1}$ with the supremum norm
    $\norm{f}_\infty = \max_{t \in [0,1]} \abs{f(t)}$, or with the $p$-norm
    $\norm{f}_p = (\int_0^1 \abs{f}^p)^{1/p}$.
  \item $X = \Contk{1}{0,1}$ with
    $\norm{f} = \max_t \abs{f(t)} + \max_t \abs{f'(t)}$.
  \item $\lp{p} = \set[\big]{x = (x_n) \,\big|\, \sum_{n=1}^{\infty}
    \abs{x_n}^p < \infty}$ with
    $\norm{x}_p = (\sum_{n=1}^{\infty}\abs{x_n}^p)^{1/p}$.
  \item $\lp{\infty}$, the bounded sequences, with
    $\norm{x}_\infty = \sup_n \abs{x_n}$.
\end{enumerate}
\end{example}

Of the norm axioms, only the triangle inequality for the $p$-norm is not
immediate. The next two theorems supply it.

\begin{theorem}[Hölder's inequality]\label{thm:holder}
Let $1 < p < \infty$ and let $q$ be its \emph{conjugate exponent}, defined by
$\frac{1}{p} + \frac{1}{q} = 1$. Then for all $x, y \in \K^n$,
\begin{equation}\label{eq:holder}
  \norm{x \cdot y}_1 \;=\; \sum_{j=1}^{n} \abs{x_j y_j}
  \;\le\; \Bigl(\sum_{j=1}^{n} \abs{x_j}^p\Bigr)^{1/p}
          \Bigl(\sum_{j=1}^{n} \abs{y_j}^q\Bigr)^{1/q}
  \;=\; \norm{x}_p \norm{y}_q ,
\end{equation}
where $x \cdot y = (x_1y_1, \dots, x_ny_n)$. The inequality also holds for the
endpoint pair $(p,q) = (1,\infty)$ and, by symmetry, for $(\infty, 1)$.
\end{theorem}

\begin{remark}
For $p = q = 2$ the inequality is the Cauchy--Schwarz inequality.
\end{remark}

\begin{proof}
The endpoint case is immediate:
$\sum_j \abs{x_j y_j} \le (\max_j \abs{y_j}) \sum_j \abs{x_j}
= \norm{x}_1 \norm{y}_\infty$.

Let now $1 < p < \infty$. The first step is Young's inequality: for all
$a, b \ge 0$,
\begin{equation}\label{eq:young}
  ab \le \frac{a^p}{p} + \frac{b^q}{q}.
\end{equation}
If $a = 0$ or $b = 0$ the left side vanishes and the right side is
non-negative, so assume $a, b > 0$ and put $u_1 = a^p > 0$, $u_2 = b^q > 0$,
$\alpha_1 = \frac{1}{p}$, $\alpha_2 = \frac{1}{q}$, so that
$\alpha_1 + \alpha_2 = 1$. The logarithm is concave on $(0,\infty)$, since
$(\log)''(t) = -1/t^2 < 0$, and concavity applied to the convex combination
$\alpha_1 u_1 + \alpha_2 u_2$ gives
\[
  \log\bigl(u_1^{\alpha_1} u_2^{\alpha_2}\bigr)
  = \alpha_1 \log u_1 + \alpha_2 \log u_2
  \le \log(\alpha_1 u_1 + \alpha_2 u_2).
\]
The logarithm is increasing, so
$u_1^{\alpha_1} u_2^{\alpha_2} \le \alpha_1 u_1 + \alpha_2 u_2$, which is
exactly \eqref{eq:young}, because
$u_1^{\alpha_1}u_2^{\alpha_2} = (a^p)^{1/p}(b^q)^{1/q} = ab$.

\begin{nfigure}{Concavity of the logarithm: the chord over $u_1$ and $u_2$
lies below the graph, so
$\alpha_1 \log u_1 + \alpha_2 \log u_2 \le \log(\alpha_1 u_1 + \alpha_2 u_2)$.}
\begin{tikzpicture}[scale=0.95]
  \draw[rulegrey] (0,0) -- (5.0,0);
  \draw[rulegrey] (0.25,-0.9) -- (0.25,1.6);
  \draw[accent, thick] plot[domain=0.32:4.8, samples=90] (\x, {ln(\x)});
  \draw[thmaccent, thick] (1.0,{ln(1.0)}) -- (4.2,{ln(4.2)});
  \draw[rulegrey, dashed] (1.0,{ln(1.0)}) -- (1.0,-0.15);
  \draw[rulegrey, dashed] (4.2,{ln(4.2)}) -- (4.2,0);
  \node[below] at (1.0,-0.15) {$u_1$};
  \node[below] at (4.2,0) {$u_2$};
  \node[accent] at (4.6,1.75) {$\log$};
\end{tikzpicture}
\end{nfigure}

The second step is a reduction by scaling. If $x = 0$ or $y = 0$ both sides of
\eqref{eq:holder} vanish, so assume $\norm{x}_p > 0$ and $\norm{y}_q > 0$ and
put $\tilde{x} = x / \norm{x}_p$ and $\tilde{y} = y / \norm{y}_q$, so that
$\norm{\tilde{x}}_p = \norm{\tilde{y}}_q = 1$. Both sides of
\eqref{eq:holder} are positively homogeneous of degree one in $x$ and in $y$
separately, so \eqref{eq:holder} for the pair $x, y$ is equivalent to
\eqref{eq:holder} for the normalised pair $\tilde{x}, \tilde{y}$. It therefore
suffices to treat the case
$\sum_j \abs{x_j}^p = \sum_j \abs{y_j}^q = 1$.

In that case, applying \eqref{eq:young} with $a = \abs{x_j}$ and
$b = \abs{y_j}$ and summing over $j$,
\[
  \sum_{j=1}^{n} \abs{x_j y_j}
  \le \sum_{j=1}^{n} \Bigl(\frac{\abs{x_j}^p}{p}
      + \frac{\abs{y_j}^q}{q}\Bigr)
  = \frac{1}{p} + \frac{1}{q} = 1 = \norm{x}_p \norm{y}_q . \qedhere
\]
\end{proof}

\begin{theorem}[Minkowski's inequality]\label{thm:minkowski}
Let $1 \le p < \infty$. Then $\norm{x+y}_p \le \norm{x}_p + \norm{y}_p$ for
all $x, y \in \K^n$.
\end{theorem}

\begin{proof}
For $p = 1$ the claim is the triangle inequality in $\K$ applied termwise. Let
$1 < p < \infty$ and let $q$ be the conjugate exponent, so that
$q = \frac{p}{p-1}$ and hence $(p-1)q = p$. If
$\sum_j \abs{x_j + y_j}^p = 0$ the left side vanishes and the claim is
trivial, so assume the sum is positive; being a finite sum, it is also finite.
Using $\abs{x_j+y_j} \le \abs{x_j} + \abs{y_j}$ and then Hölder
\eqref{eq:holder} on each of the two resulting sums, with exponents $q$ and
$p$,
\[
  \begin{aligned}
    \sum_j \abs{x_j+y_j}^p
    &\le \sum_j \abs{x_j+y_j}^{p-1}\abs{x_j}
       + \sum_j \abs{x_j+y_j}^{p-1}\abs{y_j} \\
    &\le \Bigl(\sum_j \abs{x_j+y_j}^{(p-1)q}\Bigr)^{1/q}
         \biggl[\Bigl(\sum_j \abs{x_j}^p\Bigr)^{1/p}
              + \Bigl(\sum_j \abs{y_j}^p\Bigr)^{1/p}\biggr] \\
    &= \Bigl(\sum_j \abs{x_j+y_j}^{p}\Bigr)^{1/q}
       \bigl(\norm{x}_p + \norm{y}_p\bigr).
  \end{aligned}
\]
Dividing by the positive finite number
$(\sum_j \abs{x_j+y_j}^p)^{1/q}$ and using $1 - \frac{1}{q} = \frac{1}{p}$,
\[
  \norm{x+y}_p = \Bigl(\sum_j \abs{x_j+y_j}^p\Bigr)^{1/p}
  \le \norm{x}_p + \norm{y}_p . \qedhere
\]
\end{proof}

% [p2]
\subsection{Infinite sequences and functions}

Both inequalities carry over from $\K^n$ to $\lp{p}$ and to $\Cont{0,1}$.

\begin{theorem}\label{thm:holder-inf}
Let $1 < p < \infty$ and $\frac{1}{p} + \frac{1}{q} = 1$. For all
$x \in \lp{p}$ and $y \in \lp{q}$,
$\norm{x \cdot y}_1 \le \norm{x}_p \norm{y}_q$, and for all
$f, g \in \Cont{0,1}$,
\[
  \norm{fg}_1 = \int_0^1 \abs{fg}
  \le \Bigl(\int_0^1 \abs{f}^p\Bigr)^{1/p}
      \Bigl(\int_0^1 \abs{g}^q\Bigr)^{1/q}
  = \norm{f}_p \norm{g}_q .
\]
\end{theorem}

\begin{theorem}\label{thm:minkowski-inf}
Let $1 \le p < \infty$. For all $x, y \in \lp{p}$,
$\norm{x+y}_p \le \norm{x}_p + \norm{y}_p$, and for all
$f, g \in \Cont{0,1}$, $\norm{f+g}_p \le \norm{f}_p + \norm{g}_p$.
\end{theorem}

\subsection{The spaces \texorpdfstring{$L_p[0,1]$}{Lp[0,1]}}

\begin{definition}[{$\Lp{p}[0,1]$}]\label{def:lp}
Let $1 \le p < \infty$. On $\Cont{0,1}$ the $p$-norm is a norm, but the space
$(\Cont{0,1}, \norm{\cdot}_p)$ is not complete. Its completion, in the sense
of Proposition~\ref{prop:completion}, is denoted $\Lp{p}[0,1]$.
\end{definition}

\begin{remark}\label{rem:lp-repr}
The space $\Lp{p}[0,1]$ has a concrete description. Its elements are
represented by those measurable functions $f \colon [0,1] \to \K$ for which
$\int_0^1 \abs{f}^p < \infty$, and two such functions represent the same
element exactly when they agree outside a set of measure zero.
\end{remark}

\begin{remark}\label{rem:lp-measure}
More generally, let $(X, \mathcal{A}, \mu)$ be a measure space and
$1 \le p < \infty$. Then
\[
  \Lp{p}(X, \mathcal{A}, \mu) = \set[\Big]{f \colon X \to \K \,\Big|\,
    f \text{ measurable}, \ \int_X \abs{f}^p \dd\mu < \infty},
  \qquad
  \norm{f}_p = \Bigl(\int_X \abs{f}^p \dd\mu\Bigr)^{1/p},
\]
and $\Lp{p}(X, \mathcal{A}, \mu)$ is a Banach space.
\end{remark}

\begin{exercise}\label{exc:lp-inclusion}
Let $1 \le p < p' < \infty$. Show that $\lp{p} \subset \lp{p'}$.
\end{exercise}

\begin{solution}
Let $x \in \lp{p}$. The series $\sum_j \abs{x_j}^p$ converges, so
$\abs{x_j} \to 0$, and there is an $N$ with $\abs{x_j} < 1$ for all $j > N$.
For such $j$ the larger exponent gives the smaller value,
$\abs{x_j}^{p'} \le \abs{x_j}^{p}$, whence
\[
  \sum_{j > N} \abs{x_j}^{p'} \le \sum_{j > N} \abs{x_j}^{p} < \infty .
\]
The remaining finitely many terms are finite, so
$\sum_j \abs{x_j}^{p'} < \infty$, that is $x \in \lp{p'}$. In particular
$\lp{1} \subset \lp{3/2} \subset \lp{2} \subset \dots \subset \lp{\infty}$.
\end{solution}

\begin{exercise}\label{exc:Lp-inclusion}
Let $1 \le p < p' < \infty$. Show that
$\Lp{p}[0,1] \supset \Lp{p'}[0,1]$, the inclusion running the other way than
in Exercise~\ref{exc:lp-inclusion}.
\end{exercise}

\begin{solution}
Work with the representatives of Remark~\ref{rem:lp-repr}, so that
$f \in \Lp{p'}[0,1]$ means $\int_0^1 \abs{f}^{p'} < \infty$. Split the domain
according to the size of $\abs{f}$:
\[
  \int_0^1 \abs{f}^p
  = \int_{\set{\abs{f} > 1}} \abs{f}^p
    + \int_{\set{\abs{f} \le 1}} \abs{f}^p
  \le \int_0^1 \abs{f}^{p'} + 1 < \infty .
\]
On the first set $\abs{f}^p \le \abs{f}^{p'}$, since $p < p'$ and
$\abs{f} > 1$ there; the second integral is at most the measure of $[0,1]$,
which is $1$. Hence $f \in \Lp{p}[0,1]$.
\end{solution}

% [p2]
\section{Equivalent norms}

\begin{definition}[Equivalent norms]\label{def:equiv}
Two norms $\norm{\cdot}_a$ and $\norm{\cdot}_b$ on a vector space $X$ are
\emph{equivalent} if there are constants $c, C > 0$ such that
$c\norm{x}_a \le \norm{x}_b \le C\norm{x}_a$ for all $x \in X$.
\end{definition}

Equivalent norms induce the same topology, hence the same convergent
sequences with the same limits, and the same Cauchy sequences. As a rule of
thumb, the norms occurring naturally on an infinite-dimensional space are not
equivalent; the finite-dimensional case is the opposite extreme.

\begin{lemma}\label{lem:independent}
Let $(X, \norm{\cdot})$ be a normed space over $\K$ and let
$x_1, \dots, x_n \in X$ be linearly independent. Then there is a $C > 0$ such
that
\[
  \Bigl\lVert \sum_{i=1}^{n} \alpha_i x_i \Bigr\rVert
  \;\ge\; C \sum_{i=1}^{n} \abs{\alpha_i}
  \qquad \text{for all } \alpha_1, \dots, \alpha_n \in \K .
\]
\end{lemma}

\begin{proof}
Both sides are positively homogeneous of degree one in $(\alpha_1, \dots,
\alpha_n)$, and both vanish when every $\alpha_i$ is zero, so it suffices to
find a $C > 0$ that works for all coefficient vectors with
$\sum_i \abs{\alpha_i} = 1$.

Suppose no such $C$ exists. Then for every $m \in \N$ there are coefficients
$\alpha^{(m)}_1, \dots, \alpha^{(m)}_n$ with $\sum_i \abs{\alpha^{(m)}_i} = 1$
and
\[
  \Bigl\lVert \sum_{i=1}^{n} \alpha^{(m)}_i x_i \Bigr\rVert < \frac{1}{m}.
\]
Each coordinate sequence $(\alpha^{(m)}_i)_m$ is bounded by $1$, so by the
Bolzano--Weierstrass theorem it has a convergent subsequence. Extracting
successively for $i = 1, \dots, n$ and passing to the last subsequence, we may
assume $\alpha^{(m)}_i \to \alpha_i$ for every $i$ simultaneously. Then
$\sum_i \abs{\alpha_i} = \lim_m \sum_i \abs{\alpha^{(m)}_i} = 1$, so not every
$\alpha_i$ is zero. The map $(\alpha_1, \dots, \alpha_n) \mapsto
\sum_i \alpha_i x_i$ is continuous, and so is the norm, hence
\[
  \Bigl\lVert \sum_{i=1}^{n} \alpha_i x_i \Bigr\rVert
  = \lim_{m \to \infty}
    \Bigl\lVert \sum_{i=1}^{n} \alpha^{(m)}_i x_i \Bigr\rVert = 0,
\]
so $\sum_i \alpha_i x_i = 0$ with coefficients not all zero. This contradicts
the independence of $x_1, \dots, x_n$.
\end{proof}

% [p3]
\begin{theorem}\label{thm:finite-dim}
If $X$ is a finite-dimensional vector space over $\K$, then any two norms on
$X$ are equivalent.
\end{theorem}

\begin{proof}
Let $x_1, \dots, x_n$ be a basis of $X$ and, for
$x = \sum_{i=1}^{n} \alpha_i x_i$, set $\norm{x}_1 = \sum_{i=1}^n
\abs{\alpha_i}$, which is a norm on $X$. Let $\norm{\cdot}$ be any norm on
$X$. By Lemma~\ref{lem:independent} there is a $C > 0$ with
$\norm{x} \ge C \norm{x}_1$ for all $x \in X$. In the other direction,
\[
  \norm{x} = \Bigl\lVert \sum_{i=1}^n \alpha_i x_i \Bigr\rVert
  \le \sum_{i=1}^n \abs{\alpha_i} \norm{x_i}
  \le \Bigl(\max_{1 \le i \le n} \norm{x_i}\Bigr) \sum_{i=1}^n
      \abs{\alpha_i}
  = c' \norm{x}_1 ,
\]
with $c' = \max_{1 \le i \le n} \norm{x_i} > 0$, the maximum being positive
because the basis vectors are non-zero. So every norm on $X$ is equivalent to
$\norm{\cdot}_1$.

Equivalence of norms is transitive: if
$c_a \norm{x}_1 \le \norm{x}_a \le C_a \norm{x}_1$ and
$c_b \norm{x}_1 \le \norm{x}_b \le C_b \norm{x}_1$ for all $x$, then
\[
  \frac{c_b}{C_a} \norm{x}_a \le \norm{x}_b
  \le \frac{C_b}{c_a} \norm{x}_a .
\]
Applying this to two arbitrary norms on $X$, each equivalent to
$\norm{\cdot}_1$ by the first part, gives the claim.
\end{proof}
